\documentclass[aspectratio=169,11pt]{beamer}
\usepackage{paperdeck}
\title[BERTScore]{BERTScore: Evaluating Text Generation with BERT}
\subtitle{Zhang, Kishore, Wu, Weinberger and Artzi (ICLR 2020), arXiv:1904.09675}
\author[Paper summary]{Paper summary for the CS-781 capstone}
\date{}
\begin{document}

\begin{frame}\titlepage\end{frame}

\paperglance
{T.\ Zhang, V.\ Kishore, F.\ Wu, K.\ Q.\ Weinberger, Y.\ Artzi. \emph{BERTScore: Evaluating Text Generation with BERT}. ICLR 2020. Cornell University and ASAPP.}
{Score generated text by matching each token to its most similar token in the reference, using cosine similarity of \emph{contextual embeddings} instead of exact word overlap.}
{BERTScore F1 is the \textbf{primary metric} for Phase 1 (continuations) and Phase 6 (summaries). Our \code{06\_evaluate.py} computes it locally.}
{Main deck: ``Step 06: the metrics''. Glossary: ``BERTScore'' slide.}

\begin{frame}{Roadmap}
  \begin{multicols}{2}\small\tableofcontents\end{multicols}
\end{frame}

% ------------------------------------------------------------------
\section{The problem}

\begin{frame}{What is wrong with n-gram metrics?}
  \begin{columns}[T]
    \begin{column}{0.52\textwidth}
      \small
      BLEU, ROUGE and METEOR count overlapping words or n-grams between a candidate and a reference. Two failures:
      \begin{enumerate}
        \item \textbf{Paraphrases are punished.} Correct meaning in different words scores low.
        \item \textbf{Order and long-range meaning are missed.} With short n-grams, swapping cause and effect (``A because B'' vs ``B because A'') is only mildly penalised.
      \end{enumerate}
      \vspace{1mm}
      Learned metrics (e.g.\ RUSE, BEER) fix some of this, but need human judgement data for training and may not transfer to new domains.
    \end{column}
    \begin{column}{0.44\textwidth}
      \begin{exampleblock}{The paper's example}
        \small
        Reference: \emph{people like foreign cars}.
        \begin{itemize}
          \item A: \emph{people like visiting places abroad}
          \item B: \emph{consumers prefer imported cars}
        \end{itemize}
        BLEU and METEOR score A \textbf{higher} than B, although B means the same thing and A does not.
      \end{exampleblock}
      \begin{exampleblock}{A legal version}
        \small
        Reference: \emph{the appeal is dismissed}.
        Candidate: \emph{the petition is rejected}.
        Word overlap sees only ``the'' and ``is''.
      \end{exampleblock}
    \end{column}
  \end{columns}
\end{frame}

% ------------------------------------------------------------------
\section{The basic idea}

\begin{frame}{BERTScore in three steps}
  \begin{columns}[T]
    \begin{column}{0.52\textwidth}
      \small
      \begin{enumerate}
        \item \textbf{Embed}: run reference $x$ and candidate $\hat x$ through a pretrained encoder (BERT, RoBERTa, \ldots). Each token gets a \emph{contextual} vector that depends on its sentence.
        \item \textbf{Compare}: compute the cosine similarity of every reference token with every candidate token. Vectors are pre-normalised, so this is a dot product.
        \item \textbf{Match greedily}: each token takes its best match on the other side. Average these to get recall and precision, then F1.
      \end{enumerate}
    \end{column}
    \begin{column}{0.44\textwidth}
      \centering
      \begin{tikzpicture}[node distance=3mm]
        \node[box] (r) {reference $x$};
        \node[box,right=6mm of r] (c) {candidate $\hat x$};
        \node[hot,below=of r] (er) {encoder};
        \node[hot,below=of c] (ec) {encoder};
        \node[box,below=12mm of $(r)!0.5!(c)$,minimum width=36mm] (m) {cosine matrix $x_i^\top\hat x_j$};
        \node[box,below=of m,fill=aqua!15] (s) {greedy max, then $R$, $P$, $F_1$};
        \draw[arr] (r)--(er); \draw[arr] (c)--(ec);
        \draw[arr] (er)--(m); \draw[arr] (ec)--(m); \draw[arr] (m)--(s);
      \end{tikzpicture}
    \end{column}
  \end{columns}
\end{frame}

\begin{frame}{The formulas}
  \begin{columns}[T]
    \begin{column}{0.52\textwidth}
      \small
      \[
        R_{\text{BERT}}=\frac{1}{|x|}\sum_{x_i\in x}\max_{\hat x_j\in\hat x}x_i^\top\hat x_j
      \]
      \[
        P_{\text{BERT}}=\frac{1}{|\hat x|}\sum_{\hat x_j\in\hat x}\max_{x_i\in x}x_i^\top\hat x_j
      \]
      \[
        F_{\text{BERT}}=2\,\frac{P_{\text{BERT}}\cdot R_{\text{BERT}}}{P_{\text{BERT}}+R_{\text{BERT}}}
      \]
      \begin{itemize}
        \item \textbf{Recall}: is every reference token covered by something similar in the candidate?
        \item \textbf{Precision}: is every candidate token supported by something in the reference?
      \end{itemize}
    \end{column}
    \begin{column}{0.44\textwidth}
      \begin{block}{Two optional extras}
        \small
        \textbf{idf weighting}: rare words matter more.
        \[
          \text{idf}(w)=-\log\tfrac1M\textstyle\sum_{i=1}^{M}\mathbb{I}[w\in x^{(i)}]
        \]
        computed on the $M$ test references, with plus-one smoothing.
        \vspace{1mm}

        \textbf{Baseline rescaling}: raw scores cluster in a narrow range. Subtract the average score $b$ of random sentence pairs:
        \[
          \hat R=\frac{R-b}{1-b}
        \]
        This only improves readability; rankings are unchanged.
      \end{block}
    \end{column}
  \end{columns}
\end{frame}

\begin{frame}{Easy worked example}
  \begin{columns}[T]
    \begin{column}{0.5\textwidth}
      \small
      Reference: \emph{appeal dismissed with costs}\\
      Candidate: \emph{petition rejected with costs}
      \vspace{1mm}

      Illustrative cosine similarities (content words only):
      \vspace{1mm}

      \begin{tabular}{@{}lccc@{}}
        \toprule
        & petition & rejected & costs\\
        \midrule
        appeal & \textbf{0.80} & 0.30 & 0.20\\
        dismissed & 0.30 & \textbf{0.85} & 0.25\\
        costs & 0.20 & 0.25 & \textbf{1.00}\\
        \bottomrule
      \end{tabular}
      \vspace{1mm}

      {\scriptsize Numbers chosen for illustration, not computed by a model.}
    \end{column}
    \begin{column}{0.46\textwidth}
      \begin{exampleblock}{Scoring}
        \small
        Recall: row maxima $0.80,0.85,1.00$, so $R=0.883$.\\
        Precision: column maxima $0.80,0.85,1.00$, so $P=0.883$.\\
        $F_1=0.883$.
        \vspace{1mm}

        Exact unigram matching finds only ``costs'': $P=R=1/3$. BERTScore rewards the right \emph{meaning}.
      \end{exampleblock}
      \begin{alertblock}{Watch out}
        \small
        ``Allowed'' and ``dismissed'' occur in similar contexts, so a reversed outcome may lose less score than it should.
      \end{alertblock}
    \end{column}
  \end{columns}
\end{frame}

\begin{frame}{Design choices and the paper's recommendations}
  \begin{columns}[T]
    \begin{column}{0.5\textwidth}
      \small
      \begin{itemize}
        \item \textbf{12 embedding models} were compared: BERT, RoBERTa, XLNet and XLM variants.
        \item \textbf{Which layer?} Intermediate layers work better than the last one. The best layer per model was chosen on WMT16.
        \item \textbf{Greedy vs optimal matching}: greedy is simpler and performed comparably (appendix).
        \item \textbf{Tokenisation}: each model's own word pieces.
      \end{itemize}
    \end{column}
    \begin{column}{0.46\textwidth}
      \begin{exampleblock}{Recommendations}
        \small
        \begin{itemize}
          \item Use \textbf{F1}: recall, precision and F1 alternate as best, but F1 is reliably good.
          \item English: \textbf{24-layer RoBERTa-large}.
          \item Other languages: multilingual BERT, which is less stable for low-resource languages.
          \item idf weighting sometimes helps slightly and sometimes not.
        \end{itemize}
      \end{exampleblock}
      \begin{block}{Our local setup}
        \small
        \code{bert\_score} with \code{lang="en"} defaults to RoBERTa-large, matching this advice.
      \end{block}
    \end{column}
  \end{columns}
\end{frame}

% ------------------------------------------------------------------
\section{Experiments}

\begin{frame}{How a metric is evaluated}
  \begin{columns}[T]
    \begin{column}{0.52\textwidth}
      \small
      A good metric \textbf{agrees with human judgements}. Tests used:
      \begin{itemize}
        \item \textbf{System level}: Pearson correlation between each system's average metric score and its human score.
        \item \textbf{Segment level}: Kendall correlation per sentence. Much harder.
        \item \textbf{Model selection}: in 100 random systems, does the metric pick the same best system as humans (Hits@1)?
        \item \textbf{Robustness}: can it tell real paraphrases from adversarial word swaps?
      \end{itemize}
    \end{column}
    \begin{column}{0.44\textwidth}
      \begin{block}{Data}
        \small
        \begin{itemize}
          \item WMT18 translation metrics task: 149 systems, 14 language directions. WMT16--17 add 202 more.
          \item COCO 2015 image captioning: 12 systems.
          \item QQP and PAWS paraphrase pairs.
        \end{itemize}
        363 systems in total across MT and captioning.
      \end{block}
    \end{column}
  \end{columns}
\end{frame}

\begin{frame}{Machine translation: segment-level agreement}
  \begin{columns}[T]
    \begin{column}{0.56\textwidth}
      \footnotesize
      Kendall $\tau$ with human judgements on WMT18 (to-English / from-English):
      \vspace{1mm}

      \begin{tabular}{@{}lccc@{}}
        \toprule
        Pair & BLEU & YiSi-1 & $F_{\text{BERT}}$\\
        \midrule
        cs & .233/.389 & .319/.496 & \textbf{.404/.562}\\
        de & .415/.620 & .488/.691 & \textbf{.550/.728}\\
        et & .285/.414 & .351/.546 & \textbf{.397/.586}\\
        fi & .154/.355 & .231/.504 & \textbf{.296/.546}\\
        ru & .228/.330 & .300/.407 & \textbf{.353/.423}\\
        tr & .145/.261 & .234/.418 & \textbf{.292}/.399\\
        zh & .178/.311 & .211/.323 & \textbf{.264/.364}\\
        \bottomrule
      \end{tabular}
    \end{column}
    \begin{column}{0.4\textwidth}
      \begin{exampleblock}{Reading it}
        \small
        \begin{itemize}
          \item BERTScore beats BLEU in every direction at the sentence level, often by a wide margin.
          \item It even beat RUSE, a metric \emph{trained} on human judgements.
          \item At the \emph{system} level most metrics already correlate at 0.95 or more, so gains there are smaller.
        \end{itemize}
        So BERTScore is most useful for judging \textbf{individual outputs}.
      \end{exampleblock}
    \end{column}
  \end{columns}
\end{frame}

\begin{frame}{Image captioning and paraphrase robustness}
  \begin{columns}[T]
    \begin{column}{0.48\textwidth}
      \footnotesize
      \textbf{COCO 2015}, Pearson correlation with two human measures (M1 / M2):
      \vspace{1mm}

      \begin{tabular}{@{}lcc@{}}
        \toprule
        Metric & M1 & M2\\
        \midrule
        BLEU & $-0.019$ & $-0.005$\\
        ROUGE-L & 0.090 & 0.096\\
        CIDEr & 0.438 & 0.440\\
        SPICE (task-specific) & 0.759 & 0.750\\
        $R_{\text{BERT}}$ & 0.888 & 0.863\\
        $R_{\text{BERT}}$ + idf & \textbf{0.917} & \textbf{0.889}\\
        \bottomrule
      \end{tabular}
      \vspace{1mm}

      Only LEIC, which is trained on COCO and sees the images, scored higher.
    \end{column}
    \begin{column}{0.48\textwidth}
      \footnotesize
      \textbf{Paraphrase detection}, AUC (QQP / adversarial PAWS):
      \vspace{1mm}

      \begin{tabular}{@{}lcc@{}}
        \toprule
        Method & QQP & PAWS\\
        \midrule
        BERT trained on QQP & 0.963 & 0.351\\
        BLEU & 0.707 & 0.527\\
        METEOR & 0.755 & 0.532\\
        $F_{\text{BERT}}$ & 0.761 & 0.685\\
        $F_{\text{BERT}}$ + idf & \textbf{0.777} & \textbf{0.693}\\
        \bottomrule
      \end{tabular}
      \vspace{1mm}

      PAWS swaps words: ``flights from New York to Florida'' vs ``from Florida to New York''. n-gram metrics fall to near chance (0.5); BERTScore drops much less.
    \end{column}
  \end{columns}
\end{frame}

\begin{frame}{Speed}
  \begin{columns}[T]
    \begin{column}{0.52\textwidth}
      \small
      Despite running a large model, BERTScore is practical:
      \begin{itemize}
        \item About \textbf{192.5 candidate--reference pairs per second} on one GTX-1080Ti GPU.
        \item The WMT18 EN--DE test set (2,998 sentences) takes \textbf{15.6 s}, against 5.4 s for SacreBLEU.
      \end{itemize}
      The authors conclude it is cheap enough to use during validation, e.g.\ for early stopping, as well as for testing.
    \end{column}
    \begin{column}{0.44\textwidth}
      \begin{exampleblock}{For our dev set}
        \small
        200 dev continuations of 128 words each are a few hundred pairs.
        Scoring them takes seconds on a GPU. The one-time download of RoBERTa-large (about 1.4 GB) costs more than the scoring itself.
      \end{exampleblock}
    \end{column}
  \end{columns}
\end{frame}

% ------------------------------------------------------------------
\section{Assessment and outlook}

\begin{frame}{What this means for the capstone}
  \begin{columns}[T]
    \begin{column}{0.5\textwidth}
      \small
      \begin{itemize}
        \item Phase 1 and Phase 6 are ranked by BERTScore F1, so a continuation that \textbf{paraphrases} the hidden text well can score high even with few exact word matches.
        \item Content words, such as parties, statutes and outcomes, count more than style.
        \item Decoding choices matter: very random sampling adds off-topic tokens that lower precision; very repetitive text lowers recall.
      \end{itemize}
    \end{column}
    \begin{column}{0.46\textwidth}
      \begin{block}{Experiments to run}
        \small
        \begin{itemize}
          \item Compare dev BERTScore with dev loss: do they rank checkpoints the same way?
          \item Temperature and top-$p$ sweeps, scored by BERTScore.
          \item Report $P$, $R$ and $F_1$ separately to see \emph{why} a model scores as it does.
          \item Check a few outputs with reversed outcomes to see how much BERTScore penalises them.
        \end{itemize}
      \end{block}
    \end{column}
  \end{columns}
\end{frame}

\begin{frame}{Strengths and limitations}
  \begin{columns}[T]
    \begin{column}{0.48\textwidth}
      \begin{exampleblock}{Strengths}
        \small
        \begin{itemize}
          \item Rewards paraphrases and meaning, not surface form.
          \item Task-agnostic; needs no human-judgement training data.
          \item Much better sentence-level agreement with humans than BLEU.
          \item More robust to adversarial word swaps.
          \item Available in 104 languages via multilingual BERT.
          \item Fast enough for routine use.
        \end{itemize}
      \end{exampleblock}
    \end{column}
    \begin{column}{0.48\textwidth}
      \begin{alertblock}{Limitations}
        \small
        \begin{itemize}
          \item Scores depend on the model, layer and version, so they are not comparable across setups.
          \item No single configuration wins everywhere.
          \item Greedy token matching ignores global structure.
          \item Unstable for some low-resource settings: recall for Turkish$\to$English correlated only 0.054 at system level.
          \item Measures similarity, not \textbf{factual correctness}: a fluent wrong outcome may still score well.
        \end{itemize}
      \end{alertblock}
    \end{column}
  \end{columns}
\end{frame}

\begin{frame}{Future research directions (1): from the paper and after}
  \small
  \begin{tabularx}{\textwidth}{@{}lX@{}}
    \toprule
    Direction & Details\\
    \midrule
    Task-specific metrics built on BERTScore \emph{(paper)} & Use BERTScore as a subroutine, adapted to each task's needs\\
    BERTScore as a training loss \emph{(paper)} & It is differentiable, so it can reduce the gap between the training objective and evaluation\\
    When does idf weighting help? \emph{(paper)} & Likely depends on domain and test data; left open\\
    Optimal-transport matching & MoverScore (Zhao et al., 2019) uses earth mover's distance over contextual embeddings\\
    Learned metrics on pretrained models & BLEURT (Sellam et al., 2020) and COMET (Rei et al., 2020) are fine-tuned on human ratings\\
    Generation-based and LLM judges & BARTScore (Yuan et al., 2021); LLM-as-judge methods such as G-Eval (Liu et al., 2023)\\
    Domain-specific embeddings & The paper notes SciBERT-based BERTScore for scientific summaries (Gabriel et al., 2019)\\
    \bottomrule
  \end{tabularx}
\end{frame}

\begin{frame}{Future research directions (2): open questions}
  \begin{columns}[T]
    \begin{column}{0.48\textwidth}
      \begin{block}{Open in general}
        \small
        \begin{itemize}
          \item Detecting \textbf{factual errors and negation} with embedding metrics.
          \item Evaluating long outputs: whole summaries, not sentences.
          \item Calibrated, comparable scores across models and versions.
          \item Metrics that resist being gamed when used as training targets.
        \end{itemize}
      \end{block}
    \end{column}
    \begin{column}{0.48\textwidth}
      \begin{exampleblock}{Open for legal NLP}
        \small
        \begin{itemize}
          \item Would a \textbf{legal-domain encoder} give BERTScores that agree better with lawyers' judgements?
          \item Weighting legally critical tokens: outcome, statute, party, date.
          \item Pairing BERTScore with an outcome-consistency check for summaries and continuations.
          \item Building a small human-rated legal set to measure how well BERTScore agrees with experts.
        \end{itemize}
      \end{exampleblock}
    \end{column}
  \end{columns}
\end{frame}

\begin{frame}{Summary}
  \begin{itemize}
    \item n-gram metrics punish paraphrases and miss long-range meaning.
    \item \textbf{BERTScore}: contextual embeddings, a cosine-similarity matrix, greedy max matching, then $R$, $P$ and $F_1$. Optional idf weighting and baseline rescaling.
    \item It agrees with humans much better than BLEU at the sentence level (e.g.\ EN--DE $\tau$ .550 vs .415), and ranks near the top at the system level.
    \item Strong on captioning (0.917 vs SPICE 0.759) and robust to adversarial paraphrases.
    \item Recommendation: F1 with RoBERTa-large for English.
    \item For us: the metric that ranks Phase 1 and Phase 6, so tune decoding against it on the dev set.
  \end{itemize}
\end{frame}

\begin{frame}{Reference}
  \small
  Zhang, T., Kishore, V., Wu, F., Weinberger, K.\ Q., and Artzi, Y.\ (2020). BERTScore: Evaluating Text Generation with BERT. In \emph{International Conference on Learning Representations (ICLR)}. arXiv:1904.09675.
  \vspace{3mm}

  {\footnotesize Other work cited: Papineni et al.\ (2002) BLEU; Lin (2004) ROUGE; Banerjee and Lavie (2005) METEOR; Zhao et al.\ (2019) MoverScore; Gabriel et al.\ (2019); Sellam et al.\ (2020) BLEURT; Rei et al.\ (2020) COMET; Yuan et al.\ (2021) BARTScore; Liu et al.\ (2023) G-Eval.}
\end{frame}

\end{document}
